\documentclass{report}
\usepackage{amsmath,amssymb}
\setlength{\parindent}{0mm}
\setlength{\parskip}{1em}
\begin{document}
\begin{center}
\rule{6in}{1pt} \
{\large Wim I Vanroose \\
{\bf An efficient Multigrid calculation of the Far Field Map for Schr\"odinger Equations}}

Dept Mathematics and Computer Science \\ Middelheimlaan 1 \\ 2020 Antwerpen \\ Belgium
\\
{\tt wim.vanroose@ua.ac.be}\\
Siegfried Cools\\
Bram Reps\end{center}


Predicting the outcome of collisions between small atomic and
molecular systems is of fundamental importance for many areas of
science. Understanding their dynamics is crucial for example for
plasma physics, combustion and electron driven chemistry, amongst
other examples. However, it is computationally very challenging to
predict from first principles the outcome of these collisions since it
requires the solution of high--dimensional Schr\"odinger
equation. These are known to be hard to solve even with the best
current iterative methods. To calculate, for example, the breakup
reaction rates of the hydrogen molecule requires the solution of a
7--dimensional scattering problem. The reaction rates, also known as
the cross sections, are the far field amplitudes of solution. The
development of efficient solvers for this problem remains important
challenge. In this talk we discuss the development of solvers for
these scattering based on multigrid.


The Schr\"odinger equation in $d$-dimensions is
\begin{equation} \label{eq:schr}
(-\frac{1}{2}\Delta + V(\mathbf{x}) - E)\psi(\mathbf{x}) =
\phi(\mathbf{x}), \quad \text{for} \quad \mathbf{x} \in \mathbb{R}^d
\end{equation}
where $V(\mathbf{x})$ is the potential, $\psi$ is the wave function
and $\phi$ is the right hand side that is related to the initial state
of the system.

Scattering problems are described by solutions of \eqref{eq:schr} with
a positive energy $E$ and for these energies the equation is
equivalent to a Helmholtz equation
\begin{equation}
(-\Delta -k^2(\mathbf{x}))u = f
\end{equation}
with a wave number $k^2(\mathbf{x}) = 2\left(E-V(\mathbf{x})\right)$
solved with absorbing boundary conditions.

In this talk we discusses the current practices in solving the
Schrodinger equation \cite{c2,c3} and relate them to the challenges of
solving the Helmholtz equation. We also illustrate how advances such
as Complex Shifted Laplacian \cite{c5} and complex countour multigrid
\cite{c1} are applied to solve the Schr\"odinger equation.


\begin{thebibliography}{99}
\bibitem{c2}
T.N.~Rescigno, M.~Baertschy, W.A.~Isaacs and C.W.~McCurdy,
Science, 286, p 2474 (1999).
\bibitem{c3}
W.~Vanroose, F.~Martin, T.N.~Rescigno, Science, 310, p~1787, (2005).
\bibitem{c5}
Y.A~Erlangga, C.W.~Oosterlee and C.~Vuik, SIAM journal on Scientific
Computing, 27 p 1471 (2006).
\bibitem{c1}
S.~Cools, B.~Reps and W.~Vanroose, arxiv:1211.4461 (2012).
\end{thebibliography}


\end{document}
